%----------
%	SOCIOECONOMIC IMPACT
%----------	

Esta sección explora cómo el presente TFM contribuye a diversas esferas de la sociedad y la economía, subrayando tanto los beneficios potenciales de su aplicación como los retos inherentes a la automatización.

En el plano social, fomentar narrativas periodísticas libres de sesgos de género contribuye directamente a la construcción de un imaginario colectivo más justo e igualitario. Al facilitar la detección de patrones de desigualdad, esta investigación ayuda a los medios a ofrecer una visión más equilibrada de la realidad. Este cambio positivo no solo empodera a las mujeres ofreciendo referentes diversos, sino que también mejora la calidad democrática de la información, fomentando una audiencia más crítica y mejor informada.

Desde una perspectiva económica e industrial, la herramienta propuesta resulta de gran valor para los medios de comunicación y las corporaciones públicas. Tradicionalmente, las auditorías de calidad y perspectiva de género exigen un ingente esfuerzo manual. Al automatizar este primer filtro mediante agentes especializados y LLM, las redacciones pueden procesar grandes volúmenes de texto de manera eficiente, optimizando los flujos de trabajo editoriales y reduciendo significativamente los costes operativos.

Más allá del entorno periodístico y académico, los resultados de este trabajo pueden ser utilizados por instituciones públicas y observatorios de igualdad para diseñar políticas, normativas e iniciativas orientadas a garantizar la igualdad en los medios de comunicación.

El estudio exhaustivo del presupuesto económico empleado para la viabilidad y desarrollo de este proyecto se encuentra detallado en el Anexo \ref{appendix_budget}.

\newpage